\chapter*{前言}
\addcontentsline{toc}{chapter}{前言}

本讲义对应研究生专业核心课\emph{量子多体理论}（B113C004）。课程面向应用物理、理论物理等相关专业，预修课程为量子力学、统计物理与固体物理。课内学时 32，学分 2。

量子多体问题是凝聚态物理的核心之一。从二次量子化语言出发，经过格林函数与微扰展开，再到平均场、密度泛函与线性响应，构成一套处理相互作用粒子体系的标准工具箱。本讲义在课程大纲主线之外，并入施均仁讲义的相干态与泛函积分、Giuliani--Vignale 的电子液、Mahan 的 Matsubara/图技术/电声，以及 Gross--Runge--Heinonen《Many-Particle Theory》中的全同粒子与平移不变哈密顿量、密度算符、声子与声子介导相互作用、绘景与 Gell-Mann--Low、Rayleigh--Schr\"odinger/Goldstone 能量规则、粒子--空穴、GMB、骨架自洽、$W$/穿衣顶点、准粒子谱、$G_2$/Bethe--Salpeter、有限温微扰、费米液体平衡/输运、UHF 对称性困境与 Kondo。

结构上与 \texttt{lecture\_new} 一致：每章独立目录，子节分文件编写，便于分章修订与按章编译。正文尽量使用中文叙述，关键术语在首次出现时给出英文，公式与记号遵循常见多体理论习惯。

主要参考：
\begin{itemize}
  \item Fabrizio, \textit{A course in quantum many-body theory}, Springer, 2022（课程教材）；
  \item Giuliani--Vignale, \textit{Quantum Theory of the Electron Liquid}, CUP, 2005（电子液）；
  \item Mahan, \textit{Many-Particle Physics}, Springer, 2000（Matsubara 求和、图规则、电声与 Kubo 电导）；
  \item Gross, Runge and Heinonen, \textit{Many-Particle Theory}, Adam Hilger, 1991（密度/声子、Goldstone、骨架自洽、费米液体与 Kondo 等）；
  \item 施均仁，\textit{Quantum Theory of Many Particle Systems}（北京大学讲义，2023；本地副本见 \texttt{refs/Shi\_Junren\_QTMPS.pdf}）；
  \item 以及 Negele--Orland、Fetter--Walecka、Altland--Simons 等。
\end{itemize}

教学目的是：加深学生对多体物理基本规律的理解，提高用多体理论分析具体问题的能力，帮助学生建立凝聚态物理的直觉与图像，并激发参与科学研究的兴趣。学习中尤应树立严谨规范的科学态度，用发展的眼光看待理论的演变。
